\section{Weaker aggregation: a further direction}
\label{sec:aggregation}

Exact integer counting supplies more information than existential
projection needs. However, the earlier compiler already works over
Boolean OR and AND. Changing its integer entries to bits alone changes
the polynomial overhead, not the width exponent. A useful alternative
aggregate must enable a different computation. The following standard
polynomial-identity-testing route makes that target precise.

Discard unused inputs and write $a,b$ for the used ordinary and witness
counts, with $b\ge1$. Over a field $F$ of characteristic two, define
\[
 P_x(Z_1,\ldots,Z_b)
   =\sum_{y\in\B^b} f(x,y)\prod_{i:y_i=1}Z_i.
\]
Distinct witnesses contribute distinct monomials, each with coefficient
one. Hence $P_x$ is the zero polynomial exactly when $g(x)=0$, and
its total degree is at most $b$. In contrast, unweighted parity can
vanish on a nonempty fiber: the two witnesses $00,01$ contribute
$1+Z_2$, which vanishes at $Z_2=1$ but is not the zero polynomial.

\begin{proposition}[Fingerprint transfer to an exact circuit]
\label{prop:fingerprint}
Let $F=\mathbb F_{2^\lambda}$ have size $Q>b$. Suppose a Boolean
circuit of size $T$ evaluates the $\lambda$-bit field value $P_x(z)$
from $x$ and $z\in F^b$. If an integer $t\ge1$ satisfies
\[
 2^a(b/Q)^t<1,
\]
then
\[
 \C(g)\le t(T+\lambda)-1.
\]
In particular, $Q\ge2b$ permits $t=a+1$. If a division-free
arithmetic circuit of $A$ additions and multiplications evaluates
$P_x(z)$, then $T=O((A+1)\lambda^2)$ suffices.
\end{proposition}
\begin{proof}
The polynomial zero bound gives
$\Pr_z[P_x(z)=0]\le b/Q$ whenever $P_x$ is nonzero
\citep{Schwartz1980}. For completeness, a nonzero polynomial of
total degree $d$ has this bound $d/Q$: view it as a univariate
polynomial of degree $j$ in the last variable. Its leading coefficient
has degree at most $d-j$; by induction it vanishes with probability
at most $(d-j)/Q$. Otherwise the univariate polynomial has at most
$j$ roots. The union bound proves the claim.

Choose $t$ independent evaluation points. For any positive input, the
probability that all evaluations vanish is at most $(b/Q)^t$.
A union bound over at most $2^a$ inputs proves that some fixed list
works for every $x$. Hardwire that list. OR the $\lambda$ output
bits of each evaluation and then OR the $t$ results, costing
$tT+t(\lambda-1)+(t-1)$ gates. A negative input always gives zero,
so the resulting circuit is exact on both kinds of input.

In a fixed polynomial basis for $F$, addition costs $O(\lambda)$
binary XOR gates. Schoolbook polynomial multiplication followed by
reduction modulo a fixed irreducible polynomial costs
$O(\lambda^2)$ gates. Field constants are fixed bit strings.
This proves the arithmetic-circuit conversion.
\end{proof}

The fixed list depends on the verifier; finding it is not part of
the nonuniform conclusion. For a given input, random evaluations
instead give a uniform one-sided error test when evaluation is
uniformly efficient. For $Q\ge2b$, the same argument uses $t=h+1$ when only the
$h$ formal independent signals enter the evaluator, with their $p$
preprocessing gates shared. Thus a faster field evaluator would
transfer to a faster exact projection circuit with only polynomial
overhead in the used-input counts.

\paragraph{The tags preserve the cycle budget.}
The consistency construction evaluates $P_x$ if it is interpreted
over $F$ and each witness variable receives a unary weight equal to
$1$ at zero and $Z_i$ at one. Its gate-value and equality extensions
remain unique. Attaching a unary weight tensor adds one vertex and
one edge, preserving the cycle rank. Equality splitting and the exact
local contractions remain valid over $F$. Consequently weighting
does not worsen the cycle exponent. Reusing the existing contraction
order merely recovers that exponent; an improvement requires an
algebraic evaluation method that does less work.

This differs from a generic implementation of witness isolation by
dense affine constraints, as in the classical randomized isolation
approach \citep{ValiantVazirani1986}. Directly adding $k$ dense parity
checks can cost $O(kb)$ gates. Any fine-grained use of that construction
must account for the actual added gates and consistency constraints.
Unary tags impose no such increase in cycle rank.

\begin{proposition}[No short universal linear fingerprint]
\label{prop:universal-fingerprint}
Fix a map $w:\B^b\to\mathbb F_2^L$, independently of a relation.
If $\sum_{y\in W}w(y)\ne0$ for every nonempty
$W\subseteq\B^b$, then $L\ge2^b$.
\end{proposition}
\begin{proof}
The $2^b$ vectors $w(y)$ must be linearly independent over
$\mathbb F_2$. Otherwise a nonzero dependence, whose coefficients
are zeros and ones, supplies a nonempty zero-sum subset.
\end{proof}

Concatenating several field evaluations is also such a binary linear
fingerprint. The preceding transfer avoids this obstruction because
its list is chosen for one verifier's at most $2^a$ fibers, not for
all subsets of the witness universe. This is a limitation of a fixed
universal linear aggregate, not a circuit lower bound or a lower
bound for a restricted class of small verifiers.

\paragraph{A classical example where the weaker aggregate changes the algorithm.}
Let $X=(x_{ij})$ encode a bipartite graph on $r+r$ vertices, and let
the witness be a permutation $\sigma$ certifying a perfect matching.
The verifier checks distinct column choices and
$\AND_i x_{i,\sigma(i)}$. Binary column choices give
$b=O(r\log r)$ witness bits and an $O(r^2\log r)$-gate verifier:
compare every pair of columns and use a multiplexer in each row.
Out-of-range column codes are rejected.

Assign a separate tag $Z_{ij}$ to each potential edge. In
characteristic two,
\[
 \det(x_{ij}Z_{ij})
   =\sum_{\sigma\in S_r}\prod_i x_{i,\sigma(i)}Z_{i,\sigma(i)}.
\]
All permutation signs equal one. Distinct matchings give distinct
monomials, so the polynomial is nonzero exactly when a perfect
matching exists. Its degree is $r$. Determinants over a finite field
have polynomial-size Boolean evaluation circuits: perform Gaussian
elimination with explicit zero-pivot branches and compute nonzero
inverses by repeated squaring. All loop lengths are polynomial in
$r$ and $\log Q$. Applying the proof of Proposition~\ref{prop:fingerprint}
with degree $r$, $a=r^2$, $Q\ge2r$, and $t=r^2+1$ gives an exact
polynomial-size projection circuit.

This is the classical algebraic matching phenomenon
\citep{MulmuleyVaziraniVazirani1987}, not a new matching algorithm.
Exact integer counting gives the permanent, whose computation on
zero-one matrices is $\#\mathsf P$-complete \citep{Valiant1979}.
The example explains what a successful weaker aggregate can accomplish;
it supplies no unconditional separation of polynomial time from
$\#\mathsf P$. For arbitrary $B_2$ verifiers, we have not found a
field evaluator that improves the proved $1/5$ projection rate.
